\section*{Capítulo \ref{ch:g10:chemistry-of-life} --- A composição química dos seres vivos}

\begin{solution}{exo:g10:chemistry-of-life:1}
Oxigênio (cerca de 65\%), carbono (18.5\%), hidrogênio (9.5\%) e
nitrogênio (3.2\%): juntos, 96\% da massa corporal.
\end{solution}

\begin{solution}{exo:g10:chemistry-of-life:2}
Mineral: água, cloreto de sódio, fosfato de cálcio, gás carbônico (um
composto de carbono, mas fabricado fora da vida e sem esqueleto
carbono--hidrogênio). Orgânica: glicose, o triglicerídeo, o \omterm{def:g10:chemistry-of-life:families}{aminoácido}.
\end{solution}

\begin{solution}{exo:g10:chemistry-of-life:3}
Água perdida $12.0 - 0.6 = \qty{11.4}{g}$; $11.4/12.0 = 0.95$: a folha
é 95\% água.
\end{solution}

\begin{solution}{exo:g10:chemistry-of-life:4}
(a) Solução de iodo na semente moída: azul-escuro. (b) Reagente de
biureto: violeta. (c) Esmague a semente sobre papel: uma mancha
translúcida que não seca; o Sudan III a tinge de vermelho.
\end{solution}

\begin{solution}{exo:g10:chemistry-of-life:5}
\omterm{def:g10:chemistry-of-life:families}{Carboidratos}: açúcares simples como a glicose. \omterm{def:g10:chemistry-of-life:families}{Lipídios}: glicerol e
ácidos graxos. \omterm{def:g10:chemistry-of-life:families}{Proteínas}: \omterm{def:g10:chemistry-of-life:families}{aminoácidos} (vinte tipos). \omterm{def:g10:chemistry-of-life:families}{Ácidos nucleicos}:
nucleotídeos (quatro tipos).
\end{solution}

\begin{solution}{exo:g10:chemistry-of-life:6}
Carbono: $18.5/0.02 \approx 900$ vezes mais concentrado. Nitrogênio:
$3.2/0.002 = 1600$ vezes. A matéria viva não é uma amostra do que a
cerca: ela reúne alguns elementos raros por fatores de mil, o que é uma
assinatura química da vida.
\end{solution}

\begin{solution}{exo:g10:chemistry-of-life:7}
Água: $0.62 \times 70 = \qty{43.4}{kg}$. Matéria seca: $70 - 43.4 =
\qty{26.6}{kg}$. \omterm{def:g10:chemistry-of-life:families}{Proteína}: $0.17 \times 70 = \qty{11.9}{kg}$.
\end{solution}

\begin{solution}{exo:g10:chemistry-of-life:8}
\omterm{def:g10:chemistry-of-life:families}{Carboidratos} $60 \times 17 = \qty{1020}{kJ}$; \omterm{def:g10:chemistry-of-life:families}{lipídios} $12 \times 38 =
\qty{456}{kJ}$; \omterm{def:g10:chemistry-of-life:families}{proteínas} $8 \times 17 = \qty{136}{kJ}$; total
\qty{1612}{kJ}, cerca de \qty{1600}{kJ}. Os \omterm{def:g10:chemistry-of-life:families}{lipídios} fornecem $456/1612 \approx
28\%$ da energia com apenas 12\% da massa.
\end{solution}

\begin{solution}{exo:g10:chemistry-of-life:9}
Pão: azul-escuro --- a farinha é sobretudo amido. Queijo: nenhuma
mudança de cor além do próprio tom do iodo --- \omterm{def:g10:chemistry-of-life:families}{proteínas} e \omterm{def:g10:chemistry-of-life:families}{lipídios},
sem amido. Açúcar: nenhuma mudança --- a sacarose é um açúcar simples,
não amido; o iodo só detecta a longa cadeia de amido.
\end{solution}

\begin{solution}{exo:g10:chemistry-of-life:10}
A matéria seca é sobretudo orgânica: moléculas de carbono e hidrogênio
que se combinam com o oxigênio, liberando energia, gás carbônico e
água. A cinza são os \omterm{def:g10:chemistry-of-life:mineral-organic}{sais minerais}, já plenamente combinados com o
oxigênio; nada neles pode queimar.
\end{solution}

\begin{solution}{exo:g10:chemistry-of-life:11}
O reagente de Fehling reage com um grupo químico livre carregado por
moléculas isoladas de glicose; no amido as unidades de glicose estão
unidas umas às outras e esses grupos ficam trancados dentro da cadeia,
de modo que o reagente não encontra nada. O iodo, ao contrário, se
aloja na longa cadeia enrolada do amido, que moléculas isoladas de
glicose não formam. Cada teste detecta uma estrutura, não um elemento.
\end{solution}

\begin{solution}{exo:g10:chemistry-of-life:12}
O grão absorveu água em massa. A 13\% as reservas estão sólidas e as
enzimas não conseguem trabalhar: a química do grão está desligada e ele
sobrevive dormente. A água dissolve as reservas, deixa as enzimas
agirem e leva os produtos ao embrião em crescimento: a química só
recomeça quando a água é abundante.
\end{solution}

\begin{solution}{exo:g10:chemistry-of-life:13}
Sal: $9 \times 0.5 = \qty{4.5}{g}$; glicose: $1 \times 0.5 =
\qty{0.5}{g}$. A água pura diluiria o plasma; a água entraria então nas
hemácias, que inchariam e estourariam.
\end{solution}

\begin{solution}{exo:g10:chemistry-of-life:14}
Glicogênio: $500 \times 17 = \qty{8500}{kJ}$. \omterm{def:g10:chemistry-of-life:families}{Lipídio}: $10000 \times 38 =
\qty{3.8e5}{kJ}$. Para guardar \qty{3.8e5}{kJ} como glicogênio: $380000/17
\approx \qty{22.4}{kg}$ de glicogênio, mais três vezes essa
massa de água, cerca de \qty{90}{kg} ao todo --- mais do que a pessoa. A
gordura armazena aproximadamente nove vezes mais energia por quilograma
carregado, e é por isso que os animais que precisam se mover guardam as
reservas deles como \omterm{def:g10:chemistry-of-life:families}{lipídio}.
\end{solution}

\begin{solution}{exo:g10:chemistry-of-life:15}
Os átomos são de fato banais, mas (1) as proporções não são: a matéria
viva concentra o carbono mil vezes e o nitrogênio muito mais, contra
uma crosta de oxigênio e silício; (2) as moléculas não são: só os seres
vivos montam o carbono em cadeias gigantes --- \omterm{def:g10:chemistry-of-life:families}{proteínas}, amido, DNA ---
cujo tamanho não tem equivalente mineral; (3) essas moléculas carregam
informação na ordem das unidades delas, e nenhum cristal faz isso. A
vida é quimicamente especial na organização, não nos átomos.
\end{solution}

\begin{solution}{pb:g10:chemistry-of-life:1}
\textbf{1.} Água $0.13 \times 45 = \qty{5.9}{mg}$; massa seca
\qty{39.2}{mg}.

\textbf{2.} Amido $0.70 \times 39.2 \approx \qty{27.4}{mg}$; \omterm{def:g10:chemistry-of-life:families}{proteína}
$0.12 \times 39.2 \approx \qty{4.7}{mg}$; \omterm{def:g10:chemistry-of-life:families}{lipídio} $0.02 \times 39.2
\approx \qty{0.8}{mg}$.

\textbf{3.} Iodo (azul-escuro) para o amido; biureto (violeta) para a
\omterm{def:g10:chemistry-of-life:families}{proteína}; uma mancha translúcida no papel ou o Sudan III para o \omterm{def:g10:chemistry-of-life:families}{lipídio}.

\textbf{4.} A 13\% as reservas estão sólidas e as enzimas inativas; a
água precisa dissolver as reservas e deixar as reações correrem antes
que o embrião possa crescer. A germinação começa pela absorção de água,
e os 88\% da plântula são o sinal de que a química dela está ligada.

\textbf{5.} Água $0.13 \times 8 = \qty{1.04}{t}$; matéria seca
\qty{6.96}{t}, cerca de \qty{7}{t}.

\textbf{6.} Massa $500 + 320 + 10 - 130 = \qty{700}{g}$. Água no pão:
da farinha $0.13 \times 500 = \qty{65}{g}$, mais \qty{320}{g}, menos
\qty{130}{g}: \qty{255}{g}, ou seja $255/700 \approx 36\%$.

\textbf{7.} Matéria seca da farinha $0.87 \times 500 = \qty{435}{g}$;
amido $0.70 \times 435 \approx \qty{305}{g}$; \omterm{def:g10:chemistry-of-life:families}{proteína} $0.12 \times 435 \approx
\qty{52}{g}$ (\omterm{def:g10:chemistry-of-life:families}{lipídio} cerca de \qty{9}{g}).

\textbf{8.} $305 \times 17 + 52 \times 17 + 9 \times 38 \approx 5185 +
884 + 342 \approx \qty{6400}{kJ}$.

\textbf{9.} $120/700 \approx 0.17$ do pão: cerca de \qty{1100}{kJ}, ou
seja 11\% de \qty{10000}{kJ}.

\textbf{10.} O sal é mineral: ele já está plenamente oxidado e não tem
ligações carbono--hidrogênio para queimar. No corpo, os íons dele fixam
a salinidade do sangue e são necessários aos sinais nervosos e musculares.

\textbf{11.} $0.60 \times 60 = \qty{36}{kg}$ de água.

\textbf{12.} \omterm{def:g10:chemistry-of-life:families}{Proteína} $0.17 \times 60 = \qty{10.2}{kg}$; \omterm{def:g10:chemistry-of-life:families}{lipídio} $0.16
\times 60 = \qty{9.6}{kg}$; minerais \qty{3.6}{kg}; \omterm{def:g10:chemistry-of-life:families}{carboidratos}
\qty{0.6}{kg}.

\textbf{13.} $1.2/36 \approx 3.3\%$ da água dela. O volume de sangue e
o controle da temperatura dependem dela; uma perda de 10\% é perigosa,
de modo que beber precisa repô-la em algumas horas.

\textbf{14.} Plasma $0.55 \times 4.5 \approx \qty{2.5}{L}$; água $0.9
\times 2.5 \approx \qty{2.2}{kg}$, cerca de 6\% da água do corpo. O
resto está dentro das células (cerca de dois terços de toda a água do
corpo) e no líquido entre as células.

\textbf{15.} Um bebê troca uma fração bem maior da água dele a cada dia
(bebida, urina, pele) do que um adulto, de modo que algumas horas de
perda sem ingestão tiram uma fração perigosa do total dele.

\textbf{16.} \omterm{def:g10:chemistry-of-life:families}{Carboidratos} $0.40 \times 0.6 = \qty{0.24}{kg}$; \omterm{def:g10:chemistry-of-life:families}{lipídios}
$0.77 \times 9.6 \approx \qty{7.4}{kg}$; \omterm{def:g10:chemistry-of-life:families}{proteínas} $0.53 \times 10.2
\approx \qty{5.4}{kg}$. Total cerca de \qty{13}{kg} de carbono.

\textbf{17.} $13/60 \approx 22\%$, contra 18.5\%. A concordância é boa
dadas as frações arredondadas das famílias e a variação do teor de
\omterm{def:g10:chemistry-of-life:families}{lipídio} entre indivíduos; um corpo mais magro dá um número menor.

\textbf{18.} $0.16 \times 10.2 \approx \qty{1.6}{kg}$ de nitrogênio.
Ele veio do nitrogênio do ar, capturado por bactérias do solo, absorvido
pelas plantas e comido ao longo da cadeia alimentar.

\textbf{19.} $0.89 \times 36 \approx \qty{32}{kg}$ de oxigênio, ou seja
$32/60 \approx 53\%$ da massa corporal. Só a água já responde pela
maior parte dos 65\%; o restante é o oxigênio dentro das moléculas orgânicas.

\textbf{20.} Uma pessoa de \qty{60}{kg} carrega cerca de \qty{13}{kg} de
carbono. Ele deixou o ar como gás carbônico pela fotossíntese das
plantas, e entrou no corpo pela alimentação e pela digestão ao longo da
cadeia alimentar.
\end{solution}
